% Block diagrams for the unit 7 slides, one tikzpicture per page.
%
% Built by make_figs.py, which compiles this file and renders each page to
% <name>.png at 300 dpi.  The names are the "% fig:" comment above each
% picture, in page order; make_figs.py checks that the count matches.
%
% By hand:
%   pdflatex -shell-escape block_figs.tex
%   pdftocairo -png -r 300 -singlefile -f N -l N block_figs.pdf <name>
%
% The styles are unit 6's (units/unit06_timing/figs/block_figs.tex): teal for
% combinational logic, violet for registers and arrival times, amber for the
% critical path.  loopopt_figs.tex is the older, article-style source these
% replace.

\documentclass[tikz,border=4pt]{standalone}

\usepackage{amsmath}
\usepackage[scaled=0.92]{helvet}
\renewcommand{\familydefault}{\sfdefault}
\usepackage{../../styles/course}

\definecolor{violet}{HTML}{57068C}
\definecolor{tint}{HTML}{EEE6F3}
\definecolor{teal}{HTML}{00857C}
\definecolor{tealtint}{HTML}{DCEFEC}
\definecolor{amber}{HTML}{B4530F}
\definecolor{ink}{HTML}{262626}
\definecolor{muted}{HTML}{6E6E6E}

\tikzset{
  >=stealth,
  every picture/.style={line width=0.7pt, ink},
  blk/.style={draw=teal, fill=tealtint, text=ink, minimum height=0.9cm,
              minimum width=1.1cm, align=center, font=\small},
  reg/.style={shape=register, draw=violet, fill=tint, minimum width=0.7cm,
              minimum height=1.0cm},
  w/.style={->},
  crit/.style={->, line width=2.2pt, amber, every node/.append style={text=ink}},
  critl/.style={-, line width=2.2pt, amber, every node/.append style={text=ink}},
  arr/.style={font=\small\bfseries, text=violet},
  sig/.style={font=\small\ttfamily},
  dly/.style={font=\scriptsize, text=muted},
  dot/.style={circle, fill=ink, inner sep=0pt, minimum size=3.5pt},
  stage/.style={dashed, muted, line width=0.8pt},
  stagelbl/.style={font=\small\bfseries, text=muted},
  regloc/.style={circle, draw=violet, line width=1.4pt, minimum size=7mm, inner sep=0pt},
}

\begin{document}

% ----------------------------------------------------------------------
% A series of three functions, in one cycle and as a three-stage pipeline.
% ----------------------------------------------------------------------
% fig: series_comb
\begin{tikzpicture}
  \node[sig] (x) at (-1.3,0) {x};
  \node[reg] (r0) at (0,0) {};
  \node[blk] (f1) at (2.0,0) {$f_1$};
  \node[blk] (f2) at (4.2,0) {$f_2$};
  \node[blk] (f3) at (6.4,0) {$f_3$};
  \node[sig] (y) at (8.2,0) {y};
  \draw[w] (x) -- (r0.input);
  \draw[w] (r0.output) -- node[sig, above] {x\_s0} (f1);
  \draw[w] (f1) -- node[sig, above] {z1} (f2);
  \draw[w] (f2) -- node[sig, above] {z2} (f3);
  \draw[w] (f3) -- (y);
\end{tikzpicture}

% fig: series_pipe
\begin{tikzpicture}
  \node[sig] (x) at (-1.3,0) {x};
  \node[reg] (r0) at (0,0) {};
  \node[blk] (f1) at (1.9,0) {$f_1$};
  \node[reg] (r1) at (3.7,0) {};
  \node[blk] (f2) at (5.6,0) {$f_2$};
  \node[reg] (r2) at (7.4,0) {};
  \node[blk] (f3) at (9.3,0) {$f_3$};
  \node[sig] (y) at (11.0,0) {y};
  \draw[w] (x) -- (r0.input);
  \draw[w] (r0.output) -- node[sig, above] {x\_s0} (f1);
  \draw[w] (f1) -- (r1.input);
  \draw[w] (r1.output) -- node[sig, above] {z1\_s1} (f2);
  \draw[w] (f2) -- (r2.input);
  \draw[w] (r2.output) -- node[sig, above] {z2\_s2} (f3);
  \draw[w] (f3) -- (y);
  \node[stagelbl] at (0,1.0) {s0};
  \node[stagelbl] at (3.7,1.0) {s1};
  \node[stagelbl] at (7.4,1.0) {s2};
\end{tikzpicture}

% ----------------------------------------------------------------------
% Multi-path pipelining: z1 = f1(x1), z2 = f2(x2), z3 = f3(z1, x1),
% y = f4(z2, z3).  #1: 0 the circuit, 1 with stage lines, 2 with the
% register locations circled.
% ----------------------------------------------------------------------
\newcommand{\multipath}[1]{%
\begin{tikzpicture}
  \node[sig] (x1) at (-1.3,0) {x1};
  \node[reg] (r1) at (0,0) {};
  \node[sig] (x2) at (-1.3,-1.8) {x2};
  \node[reg] (r2) at (0,-1.8) {};
  \node[blk] (f1) at (2.4,0) {$f_1$};
  \node[blk] (f2) at (2.4,-1.8) {$f_2$};
  \node[blk] (f3) at (5.9,0) {$f_3$};
  \node[blk] (f4) at (9.2,-0.9) {$f_4$};
  \node[sig] (y) at (10.9,-0.9) {y};
  \draw[w] (x1) -- (r1.input);
  \draw[w] (x2) -- (r2.input);
  \draw[-] (r1.output) -- (1.2,0);
  \node[dot] at (1.2,0) {};
  \draw[w] (1.2,0) -- (f1);
  \draw[w] (1.2,0) |- (5.9,1.15) -- (f3.north);
  \draw[w] (r2.output) -- (f2);
  \draw[w] (f1) -- node[sig, above, pos=0.2] {z1} (f3);
  \draw[w] (f3.east) -- node[sig, above, pos=0.2] {z3} (7.9,0) |- ([yshift=5pt]f4.west);
  \draw[w] (f2.east) -- node[sig, above, pos=0.08] {z2} (7.9,-1.8) |- ([yshift=-5pt]f4.west);
  \draw[w] (f4) -- (y);
  \ifnum#1>0
    \foreach \xb/\lbl in {0/s0, 4.15/s1, 7.45/s2} {
      \draw[stage] (\xb,1.75) -- (\xb,-2.6);
      \node[stagelbl] at (\xb,2.0) {\lbl};
    }
    \node[stagelbl] at (2.1,-2.85) {Stage 1};
    \node[stagelbl] at (5.8,-2.85) {Stage 2};
    \node[stagelbl] at (9.4,-2.85) {Output};
  \fi
  \ifnum#1>1
    \foreach \p in {(4.15,0), (4.15,1.15), (4.15,-1.8), (7.45,0), (7.45,-1.8)}
      \node[regloc] at \p {};
  \fi
\end{tikzpicture}}

% fig: mp_orig
\multipath{0}
% fig: mp_stages
\multipath{1}
% fig: mp_regs
\multipath{2}

% fig: mp_final
\begin{tikzpicture}
  \node[sig] (x1) at (-1.3,0) {x1};
  \node[reg] (r1) at (0,0) {};
  \node[sig] (x2) at (-1.3,-1.8) {x2};
  \node[reg] (r2) at (0,-1.8) {};
  \node[blk] (f1) at (2.4,0) {$f_1$};
  \node[blk] (f2) at (2.4,-1.8) {$f_2$};
  \node[reg] (rz1) at (4.15,0) {};
  \node[reg] (rx1) at (4.15,1.15) {};
  \node[reg] (rz2) at (4.15,-1.8) {};
  \node[blk] (f3) at (5.9,0) {$f_3$};
  \node[reg] (rz3) at (7.45,0) {};
  \node[reg] (rz22) at (7.45,-1.8) {};
  \node[blk] (f4) at (9.2,-0.9) {$f_4$};
  \node[sig] (y) at (10.9,-0.9) {y};
  \draw[w] (x1) -- (r1.input);
  \draw[w] (x2) -- (r2.input);
  \draw[-] (r1.output) -- (1.2,0);
  \node[dot] at (1.2,0) {};
  \draw[w] (1.2,0) -- (f1);
  \draw[w] (1.2,0) |- (rx1.input);
  \draw[w] (r2.output) -- (f2);
  \draw[w] (f1) -- (rz1.input);
  \draw[w] (f2) -- (rz2.input);
  \draw[w] (rz1.output) -- (f3);
  \draw[w] (rx1.output) -| (f3.north);
  \draw[w] (f3) -- (rz3.input);
  \draw[w] (rz2.output) -- (rz22.input);
  \draw[w] (rz3.output) -- (7.9,0) |- ([yshift=5pt]f4.west);
  \draw[w] (rz22.output) -- (7.9,-1.8) |- ([yshift=-5pt]f4.west);
  \draw[w] (f4) -- (y);
  \node[sig, above=1pt] at (r1.north) {x1\_s0};
  \node[sig, below=1pt] at (r2.south) {x2\_s0};
  \node[sig, above=1pt] at (rx1.north) {x1\_s1};
  \node[sig, below=1pt, xshift=-4pt] at (rz1.south) {z1\_s1};
  \node[sig, below=1pt] at (rz2.south) {z2\_s1};
  \node[sig, above=1pt] at (rz3.north) {z3\_s2};
  \node[sig, below=1pt] at (rz22.south) {z2\_s2};
\end{tikzpicture}

% ----------------------------------------------------------------------
% The Two Cycle Design practice problem's solution (loopopt.xml in
% hwdesign-soln): the single-cycle circuit and the two-stage pipeline, with
% each signal's stable time.  Register outputs are stable at t = 0.
% ----------------------------------------------------------------------
\newcommand{\tcinputs}{%
  \node[sig] (x2) at (-1.4,1.5) {x2};
  \node[reg] (r2) at (0,1.5) {};
  \node[sig, above=1pt] at (r2.north) {xreg2};
  \draw[w] (x2) -- (r2.input);
  \node[sig] (x1) at (-1.4,-1.5) {x1};
  \node[reg] (r1) at (0,-1.5) {};
  \node[sig, below=1pt] at (r1.south) {xreg1};
  \draw[w] (x1) -- (r1.input);
  \node[blk] (f1) at (3.4,2.2) {$f_1$\\[-2pt] \scriptsize 3.8 ns};
  \node[blk] (f2) at (3.4,0) {$f_2$\\[-2pt] \scriptsize 2.1 ns};
  \coordinate (a1) at ([yshift=-6pt]f1.west);
  \coordinate (a2) at ([yshift=-6pt]f2.west);
  \draw[-] (r2.output) -- node[arr, below] {0} (1.2,1.5);
  \node[dot] at (1.2,1.5) {};
  \draw[w] (1.2,1.5) |- ([yshift=6pt]f1.west);
  \draw[w] (1.2,1.5) |- ([yshift=6pt]f2.west);
  \draw[-] (r1.output) -- node[arr, below] {0} (1.9,-1.5);
  \node[dot] at (1.9,-1.5) {};
  \draw[-] (1.9,-1.5) -- (1.9,0 |- a1);
  \draw[crit] (1.9,0 |- a1) -- (a1);
  \node[dot] at (1.9,0 |- a2) {};
  \draw[w] (1.9,0 |- a2) -- (a2);
}

% fig: q_tc_single
\begin{tikzpicture}
  \tcinputs
  \node[blk] (f3) at (6.4,0) {$f_3$\\[-2pt] \scriptsize 1.5 ns};
  \draw[w] (f2.east) -- node[sig, above] {z2} node[arr, below] {2.1} (f3.west);
  \draw[w] (1.9,-1.5) -| (f3.south);
  \node[blk] (f4) at (9.6,1.1) {$f_4$\\[-2pt] \scriptsize 1.2 ns};
  \draw[crit] (f1.east) -- node[pos=0.3, sig, above] {z1} node[pos=0.3, arr, below] {3.8}
    (8.4,2.2) |- ([yshift=6pt]f4.west);
  \draw[w] (f3.east) -- node[sig, above] {z3} node[arr, below] {3.6}
    (8.4,0) |- ([yshift=-6pt]f4.west);
  \draw[crit] (f4.east) -- ++(1.4,0) node[sig, right] {y} node[midway, arr, below] {5.0};
\end{tikzpicture}

% fig: q_tc_pipe
\begin{tikzpicture}
  \tcinputs
  \node[stagelbl] at (2.6,3.4) {Stage 1};
  \node[stagelbl] at (8.6,3.4) {Stage 2};
  \draw[stage] (5.4,3.6) -- (5.4,-2.6);
  \node[reg] (rz1) at (5.4,2.2) {};
  \node[sig, above=1pt, fill=white, inner sep=1pt] at (rz1.north) {z1\_s1};
  \node[reg] (rz2) at (5.4,0) {};
  \node[sig, above=1pt, fill=white, inner sep=1pt] at (rz2.north) {z2\_s1};
  \node[reg] (rx1) at (5.4,-1.5) {};
  \node[sig, below=1pt, fill=white, inner sep=1pt] at (rx1.south) {xreg1\_s1};
  \draw[crit] (f1.east) -- node[sig, above] {z1} node[arr, below] {3.8} (rz1.input);
  \draw[w] (f2.east) -- node[sig, above] {z2} node[arr, below] {2.1} (rz2.input);
  \draw[w] (1.9,-1.5) -- (rx1.input);
  \node[blk] (f3) at (8,-0.6) {$f_3$\\[-2pt] \scriptsize 1.5 ns};
  \draw[w] (rz2.output) -- node[arr, below] {0} (6.9,0) |- ([yshift=6pt]f3.west);
  \draw[w] (rx1.output) -- node[arr, below] {0} (6.9,-1.5) |- ([yshift=-6pt]f3.west);
  \node[blk] (f4) at (11,1.1) {$f_4$\\[-2pt] \scriptsize 1.2 ns};
  \draw[w] (rz1.output) -- node[pos=0.25, arr, below] {0} (9.8,2.2) |- ([yshift=6pt]f4.west);
  \draw[w] (f3.east) -- node[sig, above] {z3} node[arr, below] {1.5}
    (9.8,-0.6) |- ([yshift=-6pt]f4.west);
  \draw[w] (f4.east) -- ++(1.4,0) node[sig, right] {y} node[midway, arr, below] {2.7};
\end{tikzpicture}

% ----------------------------------------------------------------------
% The early-exit loop pipelined by hand (16-bit x): iteration i reads x[i]
% into xr, squares the xr of the iteration before into zr, and tests and
% writes the zr of the one before that.  The exit test reads only the j
% register: a compare after a register, short whatever the multiply takes.
% ----------------------------------------------------------------------
% fig: exit_pipe
\begin{tikzpicture}
  % stages, first, so the registers and labels draw over the lines
  \foreach \xb in {0, 3.8} \draw[stage] (\xb,1.9) -- (\xb,-4.4);
  \node[stagelbl] at (-1.0,2.15) {read x[i]};
  \node[stagelbl] at (1.9,2.15) {square x[i$-$1]};
  \node[stagelbl] at (7.2,2.15) {test and write x[i$-$2]};
  \node[sig] (x) at (-1.5,0) {x[i]};
  \node[reg] (rx) at (0,0) {};
  \node[sig, above=1pt, fill=white, inner sep=1pt] at (rx.north) {xr};
  \node[blk] (sq) at (1.9,0) {square};
  \node[reg] (rz) at (3.8,0) {};
  \node[sig, above=1pt, fill=white, inner sep=1pt] at (rz.north) {zr};
  \node[blk] (cmp) at (5.9,0) {$>$ zmin};
  \node[blk, minimum width=1.5cm] (wr) at (9.0,0) {write\\[-2pt] z[j]};
  \draw[w] (x) -- (rx.input);
  \draw[w] (rx.output) -- (sq);
  \draw[w] (sq) -- (rz.input);
  \draw[-] (rz.output) -- (4.6,0);
  \node[dot] at (4.6,0) {};
  \draw[w] (4.6,0) -- (cmp);
  \draw[w] (4.6,0) |- (9.0,1.1) -- node[sig, right, pos=0.5] {data} (wr.north);
  \draw[w] (cmp) -- node[sig, above] {hit} (wr);
  % the j counter: a register, enabled by hit, through +1
  \node[reg] (rj) at (5.9,-2.3) {};
  \node[sig, below=1pt] at (rj.south) {j};
  \draw[w] (cmp.south) -- node[sig, right] {en} (rj.north);
  \node[dot] (dj) at (7.2,-2.3) {};
  \draw[-] (rj.output) -- (dj);
  \draw[w] (dj) |- (8.5,-1.2) -| node[sig, right, pos=0.75] {addr} ([xshift=-12pt]wr.south);
  \node[blk] (inc) at (5.9,-3.8) {$+1$};
  \draw[w] (dj) |- (inc.east);
  \draw[w] (inc.west) -- (4.7,-3.8) |- (rj.input);
  % the exit: one compare after the j register
  \node[blk] (ex) at (9.0,-2.3) {$\geq$ NZ};
  \draw[crit] (dj) -- (ex);
  \draw[crit] (ex.east) -- ++(1.3,0) node[sig, right] {break};
\end{tikzpicture}

\end{document}
